% Generated by thesis/ch7/appendix/heart_splits_table.py -- do not edit by hand.
\begin{table}[tbp]
  \centering
  \footnotesize
  \setlength{\tabcolsep}{4pt}
  \begin{adjustbox}{max width=\textwidth}
  \begin{tabular}{@{}lccccc@{}}
    \toprule
    & Paper split & \multicolumn{4}{c}{100 random splits (242 / 61)} \\
    \cmidrule(l){3-6}
    Model & test AUC & mean & SD & 5th--95th pct. & beats $L_2$ LR \\
    \midrule
    Logistic regr., $L_2$ & 0.919 & 0.902 & 0.035 & 0.848--0.948 & -- \\
    Logistic regr., $L_1$ (sparse) & 0.926 & 0.900 & 0.037 & 0.845--0.951 & 34\% \\
    Random forest & 0.949 & 0.904 & 0.035 & 0.853--0.951 & 54\% \\
    Gradient boosting & 0.945 & 0.893 & 0.034 & 0.835--0.939 & 30\% \\
    Explainable boosting & 0.947 & 0.905 & 0.036 & 0.847--0.957 & 64\% \\
    \midrule
    PropNEAT (median, ten seeds) & 0.892 & -- & -- & -- & -- \\
    \bottomrule
  \end{tabular}
  \end{adjustbox}
  \caption{Heart: test AUC on the paper's split against 100 random stratified splits of the same 303 patients into 242 training and 61 test rows (\texttt{StratifiedShuffleSplit}, seed 1). Each comparator keeps the setting its cross-validation chose on the paper split. Across random splits every comparator averages about 0.90, and the spread between splits (SD about 0.035) is larger than any difference between models; gradient boosting has the lowest mean. The paper split favours every comparator, the tree ensembles most (by 0.04--0.05). PropNEAT was run on the paper split only. Source: \texttt{experiments/benchmarks/heart\_splits.py}.}
  \label{tab:heart-splits}
\end{table}
